%% March 2018
%%%%%%%%%%%%%%%%%%%%%%%%%%%%%%%%%%%%%%%%%%%%%%%%%%%%%%%%%%%%%%%%%%%%%%%%%%%%
% AGUJournalTemplate.tex: this template file is for articles formatted with LaTeX
%
% This file includes commands and instructions
% given in the order necessary to produce a final output that will
% satisfy AGU requirements, including customized APA reference formatting.
%
% You may copy this file and give it your
% article name, and enter your text.
%
%
% Step 1: Set the \documentclass
%
% There are two options for article format:
%
% PLEASE USE THE DRAFT OPTION TO SUBMIT YOUR PAPERS.
% The draft option produces double spaced output.
%

%% To submit your paper:
\documentclass[draft,linenumbers]{agujournal2018}
\usepackage{apacite}
\usepackage{url} %this package should fix any errors with URLs in refs.
%%%%%%%
% As of 2018 we recommend use of the TrackChanges package to mark revisions.
% The trackchanges package adds five new LaTeX commands:
%
%  \note[editor]{The note}
%  \annote[editor]{Text to annotate}{The note}
%  \add[editor]{Text to add}
%  \remove[editor]{Text to remove}
%  \change[editor]{Text to remove}{Text to add}
%
% complete documentation is here: http://trackchanges.sourceforge.net/
%%%%%%%


%% Enter journal name below.
%% Choose from this list of Journals:
%
% JGR: Atmospheres
% JGR: Biogeosciences
% JGR: Earth Surface
% JGR: Oceans
% JGR: Planets
% JGR: Solid Earth
% JGR: Space Physics
% Global Biogeochemical Cycles
% Geophysical Research Letters
% Paleoceanography and Paleoclimatology
% Radio Science
% Reviews of Geophysics
% Tectonics
% Space Weather
% Water Resources Research
% Geochemistry, Geophysics, Geosystems
% Journal of Advances in Modeling Earth Systems (JAMES)
% Earth's Future
% Earth and Space Science
% Geohealth
%
% ie, \journalname{Water Resources Research}

\journalname{JGR: Solid Earth}


% tightlist command for lists without linebreak
\providecommand{\tightlist}{%
  \setlength{\itemsep}{0pt}\setlength{\parskip}{0pt}}

% From pandoc table feature
\usepackage{longtable,booktabs,array}
\usepackage{calc} % for calculating minipage widths
% Correct order of tables after \paragraph or \subparagraph
\usepackage{etoolbox}
\makeatletter
\patchcmd\longtable{\par}{\if@noskipsec\mbox{}\fi\par}{}{}
\makeatother
% Allow footnotes in longtable head/foot
\IfFileExists{footnotehyper.sty}{\usepackage{footnotehyper}}{\usepackage{footnote}}
\makesavenoteenv{longtable}


\usepackage{amsmath}
\usepackage{amsfonts}
\usepackage{amssymb}
\providecommand{\listfigurename}{List of Figures}
\providecommand{\listtablename}{List of Tables}
\newcounter{none}
\setcounter{none}{0}
\makeatletter
\@ifpackageloaded{subcaption}{}{\usepackage{subcaption}}
\@ifpackageloaded{caption}{}{\usepackage{caption}}
\captionsetup[subfigure]{margin=0.5em}
\AtBeginDocument{%
\renewcommand*\figurename{Figure}
\renewcommand*\tablename{Table}
}
\AtBeginDocument{%
\renewcommand*\listfigurename{List of Figures}
\renewcommand*\listtablename{List of Tables}
}
\newcounter{pandoccrossref@subfigures@footnote@counter}
\newenvironment{pandoccrossrefsubfigures}{%
\setcounter{pandoccrossref@subfigures@footnote@counter}{0}
\begin{figure}\centering%
\gdef\global@pandoccrossref@subfigures@footnotes{}%
\DeclareRobustCommand{\footnote}[1]{\footnotemark%
\stepcounter{pandoccrossref@subfigures@footnote@counter}%
\ifx\global@pandoccrossref@subfigures@footnotes\empty%
\gdef\global@pandoccrossref@subfigures@footnotes{{##1}}%
\else%
\g@addto@macro\global@pandoccrossref@subfigures@footnotes{, {##1}}%
\fi}}%
{\end{figure}%
\addtocounter{footnote}{-\value{pandoccrossref@subfigures@footnote@counter}}
\@for\f:=\global@pandoccrossref@subfigures@footnotes\do{\stepcounter{footnote}\footnotetext{\f}}%
\gdef\global@pandoccrossref@subfigures@footnotes{}}
\@ifpackageloaded{float}{}{\usepackage{float}}
\floatstyle{ruled}
\@ifundefined{c@chapter}{\newfloat{codelisting}{h}{lop}}{\newfloat{codelisting}{h}{lop}[chapter]}
\floatname{codelisting}{Listing}
\newcommand*\listoflistings{\listof{codelisting}{List of Listings}}
\makeatother

\begin{document}


%% ------------------------------------------------------------------------ %%
%  Title
%
% (A title should be specific, informative, and brief. Use
% abbreviations only if they are defined in the abstract. Titles that
% start with general keywords then specific terms are optimized in
% searches)
%
%% ------------------------------------------------------------------------ %%

% Example: \title{This is a test title}

\title{Beyond Equilibrium: Kinetic Thresholds and Rheological Feedbacks
Create a Potentially Complex 410 in Slab Regions}

%% ------------------------------------------------------------------------ %%
%
%  AUTHORS AND AFFILIATIONS
%
%% ------------------------------------------------------------------------ %%

% Authors are individuals who have significantly contributed to the
% research and preparation of the article. Group authors are allowed, if
% each author in the group is separately identified in an appendix.)

% List authors by first name or initial followed by last name and
% separated by commas. Use \affil{} to number affiliations, and
% \thanks{} for author notes.
% Additional author notes should be indicated with \thanks{} (for
% example, for current addresses).

% Example: \authors{A. B. Author\affil{1}\thanks{Current address, Antartica}, B. C. Author\affil{2,3}, and D. E.
% Author\affil{3,4}\thanks{Also funded by Monsanto.}}

\authors{
Buchanan Kerswell
\affil{1}
John Wheeler
\affil{1}
Rene Gassmöller
\affil{2}
J. Huw Davies
\affil{3}
Isabel Papanagnou
\affil{4}
Sanne Cottaar
\affil{4}
}


% \affiliation{1}{First Affiliation}
% \affiliation{2}{Second Affiliation}
% \affiliation{3}{Third Affiliation}
% \affiliation{4}{Fourth Affiliation}

\affiliation{1}{Department of Earth, Oceans and Ecological Sciences,
University of Liverpool, Liverpool L69 3BX, UK}
\affiliation{2}{GEOMAR Helmholtz Centre for Ocean Research Kiel, Kiel,
Germany}
\affiliation{3}{School of Earth and Environmental Sciences, Cardiff
University, Park Place, Cardiff, CF10 3AT, UK}
\affiliation{4}{Bullard Laboratories, Department of Earth Sciences,
University of Cambridge, Cambridge CB3 0EZ, UK}
%(repeat as many times as is necessary)

%% Corresponding Author:
% Corresponding author mailing address and e-mail address:

% (include name and email addresses of the corresponding author.  More
% than one corresponding author is allowed in this LaTeX file and for
% publication; but only one corresponding author is allowed in our
% editorial system.)

% Example: \correspondingauthor{First and Last Name}{email@address.edu}
\correspondingauthor{Buchanan Kerswell}{b.kerswell@liverpool.ac.uk}

%% Keypoints, final entry on title page.

%  List up to three key points (at least one is required)
%  Key Points summarize the main points and conclusions of the article
%  Each must be 100 characters or less with no special characters or punctuation

% Example:
% \begin{keypoints}
% \item	List up to three key points (at least one is required)
% \item	Key Points summarize the main points and conclusions of the article
% \item	Each must be 100 characters or less with no special characters or punctuation
% \end{keypoints}

\begin{keypoints}
\item Plumes produce sharp 410s regardless of reaction rates; slabs show
three distinct kinetic regimes controlling discontinuity structure
\item Rheology modulates kinetic effects: strong slabs descend slowly
allowing complete reaction; weak slabs descend fast amplifying
metastability
\item Seismic observations of 410 structure in subduction zones can
constrain reaction rates but require independent rheological constraints
\end{keypoints}

%% ------------------------------------------------------------------------ %%
%
%  ABSTRACT
%
% A good abstract will begin with a short description of the problem
% being addressed, briefly describe the new data or analyses, then
% briefly states the main conclusion(s) and how they are supported and
% uncertainties.
%% ------------------------------------------------------------------------ %%

%% \begin{abstract} starts the second page

\begin{abstract}
The seismic expression of Earth's 410 km discontinuity varies across
tectonic settings, from sharp, high-amplitude interfaces to broad
transitions---patterns that cannot be explained by equilibrium
thermodynamics without invoking large-scale thermal or compositional
heterogeneities. Laboratory experiments show the olivine
\(\Leftrightarrow\) wadsleyite phase transition responsible for the 410
is rate-limited, yet previous numerical studies have not directly
evaluated the sensitivity of 410 structure to kinetic and rheological
factors. Here we investigate these relationships by coupling a
grain-scale, interface-controlled olivine \(\Leftrightarrow\) wadsleyite
growth model to compressible simulations of mantle plumes and subducting
slabs. We vary kinetic parameters across seven orders of magnitude and
quantify the resulting 410 displacements and widths. Our results reveal
an asymmetry between hot and cold environments. In plumes, high
temperatures produce sharp 410s (2--3 km wide) regardless of kinetics.
In slabs, kinetics exert first-order control on 410 structure through
three regimes: (1) quasi-equilibrium conditions producing narrow,
uplifted 410s and continuous slab descent; (2) intermediate reaction
rates generating broader, deeper 410s with metastable olivine wedges
resisting downward slab motion; and (3) ultra-sluggish reaction rates
causing slab stagnation with re-sharpened, deeply displaced 410s
(\(\lesssim\) 100 km). Rheological contrasts modulate these kinetic
effects by controlling slab geometry and residence time in the phase
transition zone. These findings demonstrate that reaction rates strongly
influence 410 structure in subduction zones, establishing the 410 as a
potential seismological constraint on upper mantle kinetic processes,
particularly in cold environments where disequilibrium effects are
amplified.
\end{abstract}
\noindent{\bf Plain language summary}\vskip-\parskip

\noindent{Seismic imaging reveals that Earth's 410 km discontinuity---a
boundary in Earth's mantle where olivine transforms into a denser form
called wadsleyite---varies significantly, sometimes appearing sharp and
other times as a broad zone. While often explained through temperature
and compositional differences, laboratory experiments show that this
mineral transformation takes time to complete, suggesting that reaction
speed (kinetics) is important. We simulated rising mantle plumes and
descending slabs to explore how kinetics and rock strength affect the
410 km discontinuity. Our models reveal striking differences. In hot
plumes, high temperatures allow fast reactions, maintaining a
consistently sharp discontinuity (2--3 km thick). In cold subducting
slabs, however, kinetics are critical. Fast reactions produce a narrow,
uplifted discontinuity as slabs sink smoothly. Moderate reaction rates
create broader, deeper discontinuities with wedges of unreacted olivine
that slow downward slab motion. Extremely slow reactions cause slabs to
stagnate with sharp 410 km discontinuities that are shifted downwards by
up to 100 km. These findings suggest that observations of the 410 km
discontinuity, particularly in subduction zones, could provide valuable
constraints on reaction rates deep within Earth's mantle.}
\vskip18pt
\clearpage

\section*{Definition of Symbols}\label{sec:symbols}
\addcontentsline{toc}{section}{Definition of Symbols}

{\def\LTcaptype{none} % do not increment counter
\begin{longtable}[]{@{}
  >{\raggedright\arraybackslash}p{(\linewidth - 6\tabcolsep) * \real{0.4688}}
  >{\raggedright\arraybackslash}p{(\linewidth - 6\tabcolsep) * \real{0.2188}}
  >{\raggedright\arraybackslash}p{(\linewidth - 6\tabcolsep) * \real{0.1562}}
  >{\raggedright\arraybackslash}p{(\linewidth - 6\tabcolsep) * \real{0.1562}}@{}}
\toprule\noalign{}
\begin{minipage}[b]{\linewidth}\raggedright
Parameter
\end{minipage} & \begin{minipage}[b]{\linewidth}\raggedright
Symbol
\end{minipage} & \begin{minipage}[b]{\linewidth}\raggedright
Unit
\end{minipage} & \begin{minipage}[b]{\linewidth}\raggedright
Equations
\end{minipage} \\
\midrule\noalign{}
\endhead
\bottomrule\noalign{}
\endlastfoot
Activation energy & \(E^{\ast}\) & J mol\(^{-1}\) &
\ref{eq:arrhenius-viscosity} \\
Activation enthalpy & \(H^{\ast}\) & J mol\(^{-1}\) &
\ref{eq:growth-rate}, \ref{eq:reaction-rate} \\
Activation factor (rheology) & \(B\) & - & \ref{eq:rheological-model} \\
Activation volume & \(V^{\ast}\) & m\(^3\) mol\(^{-1}\) &
\ref{eq:growth-rate}, \ref{eq:reaction-rate} \\
Compressibility (reference) & \(\bar{\beta}\) & Pa\(^{-1}\) &
\ref{eq:density-ala} \\
Density & \(\rho\) & kg m\(^{-3}\) &
\ref{eq:navier-stokes-no-inertia}--\ref{eq:continuity-expanded},
\ref{eq:density-ala-expansion}--\ref{eq:density-ala} \\
Density (reference) & \(\bar{\rho}\) & kg m\(^{-3}\) &
\ref{eq:adiabatic-pressure}--\ref{eq:density-ala} \\
Density (dynamic) & \(\hat{\rho}\) & kg m\(^{-3}\) & - \\
Deviatoric stress tensor & \(\sigma^{\prime}\) & Pa &
\ref{eq:navier-stokes-no-inertia}, \ref{eq:energy} \\
Deviatoric strain rate tensor & \(\dot{\epsilon}^{\prime}\) & s\(^{-1}\)
& \ref{eq:energy} \\
Gas constant & \(R\) & J mol\(^{-1}\) K\(^{-1}\) & \ref{eq:growth-rate},
\ref{eq:reaction-rate},
\ref{eq:arrhenius-viscosity}--\ref{eq:rheological-model} \\
Grain size & \(d\) & m & - \\
Gravitational acceleration & \(g\) & m s\(^{-2}\) &
\ref{eq:navier-stokes-no-inertia},
\ref{eq:adiabatic-temperature}--\ref{eq:adiabatic-pressure} \\
Growth rate & \(\dot{x}\) & m s\(^{-1}\) &
\ref{eq:volume-fraction}--\ref{eq:growth-rate},
\ref{eq:reaction-rate-short} \\
Latent heat & \(Q_L\) & J kg\(^{-1}\) & \ref{eq:energy} \\
Molar entropy & \(\bar{S}\) & J mol\(^{-1}\) K\(^{-1}\) &
\ref{eq:gibbs} \\
Molar Gibbs free energy & \(\bar{G}\) & J mol\(^{-1}\) &
\ref{eq:gibbs} \\
Molar volume & \(\bar{V}\) & m\(^{3}\) mol\(^{-1}\) & \ref{eq:gibbs} \\
Nucleation site factor & \(N\) & m\(^{-1}\) & \ref{eq:volume-fraction},
\ref{eq:reaction-rate-short} \\
Prefactor (growth rate) & \(\Gamma\) & m s\(^{-1}\) K\(^{-1}\)
ppm\(_\mathrm{OH}^{-n}\) & \ref{eq:growth-rate} \\
Prefactor (kinetic) & \(Z\) & K s\(^{-1}\) & \ref{eq:reaction-rate} \\
Prefactor (viscosity) & \(A\) & Pa s &
\ref{eq:arrhenius-viscosity}--\ref{eq:background-viscosity} \\
Pressure & \(P\) & Pa & \ref{eq:navier-stokes-no-inertia},
\ref{eq:energy}, \ref{eq:growth-rate}, \ref{eq:reaction-rate} \\
Pressure (reference) & \(\bar{P}\) & Pa & \ref{eq:adiabatic-pressure} \\
Pressure (dynamic) & \(\hat{P}\) & Pa & \ref{eq:density-ala},
\ref{eq:gibbs} \\
Reaction rate & \(\frac{dX}{dt}\), \(\frac{\partial X}{\partial t}\),
\(\dot{X}\) & s\(^{-1}\) &
\ref{eq:reaction-rate-short}--\ref{eq:composition} \\
Specific heat capacity (reference) & \(\bar{C}_p\) & J kg\(^{-1}\)
K\(^{-1}\) & \ref{eq:energy}, \ref{eq:adiabatic-temperature} \\
Temperature & \(T\) & K & \ref{eq:energy}, \ref{eq:growth-rate},
\ref{eq:reaction-rate}, \ref{eq:arrhenius-viscosity},
\ref{eq:rheological-model} \\
Temperature (reference) & \(\bar{T}\) & K &
\ref{eq:adiabatic-temperature},
\ref{eq:arrhenius-viscosity-expanded}--\ref{eq:rheological-model} \\
Temperature (dynamic) & \(\hat{T}\) & K & \ref{eq:density-ala},
\ref{eq:gibbs}, \ref{eq:arrhenius-viscosity-expanded},
\ref{eq:rheological-model} \\
Thermal conductivity (reference) & \(\bar{k}\) & W m\(^{-1}\) K\(^{-1}\)
& \ref{eq:energy} \\
Thermal expansivity (reference) & \(\bar{\alpha}\) & K\(^{-1}\) &
\ref{eq:energy}, \ref{eq:adiabatic-temperature}, \ref{eq:density-ala} \\
Time & \(t\) & s &
\ref{eq:continuity-compressible}--\ref{eq:continuity-expanded},
\ref{eq:volume-fraction},
\ref{eq:reaction-rate-short}--\ref{eq:composition} \\
Velocity & \(\vec{u}\) & m s\(^{-1}\) &
\ref{eq:continuity-compressible}--\ref{eq:continuity-expanded},
\ref{eq:composition} \\
Viscosity & \(\eta\) & Pa s &
\ref{eq:arrhenius-viscosity}--\ref{eq:arrhenius-viscosity-expanded},
\ref{eq:rheological-model} \\
Viscosity (reference) & \(\bar{\eta}\) & Pa s &
\ref{eq:background-viscosity}--\ref{eq:rheological-model} \\
Volume fraction & \(X\) & - & \ref{eq:volume-fraction},
\ref{eq:reaction-rate-short}--\ref{eq:composition} \\
Water content & \(C_\mathrm{OH}\) & ppm & \ref{eq:growth-rate} \\
Water content exponent & \(n\) & - & \ref{eq:growth-rate} \\
\end{longtable}
}

\clearpage

\section{Introduction}\label{sec:introduction}

Earth's mantle transition zone hosts two prominent seismic
discontinuities near 410 and 660 km depth, attributed to polymorphic
phase transitions of olivine \citep{ringwood1975, katsura2004}. While
these discontinuities are observed globally, their detailed seismic
characteristics---depth, sharpness, amplitude, and lateral
continuity---vary substantially between tectonic settings
\citep{deuss2009, lawrence2008, schmerr2007, fukao2013}. Some regions
display sharp, high-amplitude reflectors consistent with abrupt
mineralogical boundaries, while others exhibit broad, weakened, or
laterally variable signals. Such heterogeneity cannot be explained by
equilibrium thermodynamics alone, which relates discontinuity topography
mainly to temperature-dependent phase boundaries
\citep[e.g.,][]{cottaar2016, jenkins2016}. Additional physical
processes, including reaction kinetics and dynamic pressure effects
\citep{rubie1994, faccenda2017}, or compositional heterogeneities,
including variations in water content
\citep{karato2011, smyth1987, smyth2002} or bulk composition
\citep{glasgow2024, goes2022, saikia2008, schmerr2007, tauzin2017}
likely contribute to observed variability.

Laboratory studies demonstrate that the olivine \(\Leftrightarrow\)
wadsleyite phase transition at 410 km depth (the ``410'') occurs over
finite time scales rather than instantaneously, with kinetics governed
by temperature, pressure, water content, bulk chemical composition,
grain size, and microstructural evolution
\citep{rubie1994, liu1998, kubo2004, hosoya2005, perrillat2013, ledoux2023}.
In cold subducting slabs, sluggish reaction rates can allow metastable
olivine to persist tens of kilometers below its thermodynamic stability
limit, promoting slab stagnation and potentially triggering deep
earthquakes via transformational faulting
\citep{rubie1994, green1995, kirby1996, ishii2021, ohuchi2022, sindhusuta2025}.
In hot upwellings, slow kinetics may broaden and uplift the
discontinuity, possibly explaining reduced seismic amplitudes beneath
some hotspots \citep{chambers2005}. However, published kinetic models
remain poorly constrained, with parameters spanning several orders of
magnitude \citep[e.g.,][]{hosoya2005}, leaving the effects of
micro-scale kinetic processes on flow dynamics and seismic observables
ambiguous.

Bridging the gap between laboratory-derived kinetic rate laws and
mantle-scale seismic observations requires numerical models that couple
reaction kinetics to realistic treatments of mantle convection. Previous
modeling efforts have demonstrated that kinetics can strongly influence
mantle flow
\citep{dassler1996a, dassler1996b, schmeling1999, guest2004, agrusta2017, faccenda2017},
but investigations quantifying the sensitivity of 410 structure to
kinetic parameters remain limited. Moreover, most prior studies impose
kinematic restrictions and/or employ simplified treatments of
compressibility or kinetic rate laws that may inadequately capture
feedbacks between kinetically controlled phase transitions and flow
dynamics. Additionally, the role of rheology in modulating kinetic
effects through its control on slab geometry and descent rate has not
been thoroughly explored.

This study aims to clarify these issues by implementing a grain-scale
interface-controlled olivine \(\Leftrightarrow\) wadsleyite growth model
\citep[after][]{hosoya2005} within compressible mantle flow simulations
using the open-source geodynamic modeling software ASPECT. We
systematically explore how reaction kinetics and rheological strength
jointly influence 410 structure and address three specific questions:

\begin{enumerate}
\def\labelenumi{\arabic{enumi}.}
\tightlist
\item
  How do kinetic factors impact flow dynamics and shape the 410 in hot
  versus cold environments?
\item
  How do viscosity contrasts modulate these kinetic effects?
\item
  Can seismic observations of 410 structure constrain effective kinetic
  and rheological parameters?
\end{enumerate}

To investigate these questions, we analyze a suite of numerical
experiments varying kinetic parameters across seven orders of magnitude.
For each experiment, we test a range of viscosity contrasts and quantify
410 displacement and width, enabling direct comparisons with
seismological observations. Our results establish quantitative
relationships between reaction rates, rheological strength, flow
dynamics, and 410 structure, demonstrating that realistic treatment of
reaction kinetics is essential for accurately modeling subduction
dynamics and interpreting seismic structures.

\section{Methods}\label{sec:methods}

\subsection{Governing Equations for Compressible Mantle
Flow}\label{sec:governing-equations}

Mantle flow is simulated using the finite-element geodynamic code ASPECT
\citep[v3.0.0,][]{aspect-doi-v3.0.0, aspectmanual, heister2017, kronbichler2012, gassmoller2018, clevenger2021, fraters2019, fraters2020}
to find the velocity \(\vec{u}\), pressure \(P\), and temperature \(T\)
fields that satisfy the following equations:

\begin{equation}
  \nabla P - \nabla \cdot \sigma^{\prime} = \rho\, g
  \label{eq:navier-stokes-no-inertia}
\end{equation}

\begin{equation}
  \frac{\partial \rho}{\partial t} + \nabla \cdot (\rho\, \vec{u}) = 0
  \label{eq:continuity-compressible}
\end{equation}

\begin{equation}
  \rho\, \bar{C}_p \left(\frac{\partial T}{\partial t} + \vec{u} \cdot \nabla T \right) - \nabla \cdot \left(\bar{k}\, \nabla T \right) = \sigma^{\prime} : \dot{\epsilon}^{\prime} + \bar{\alpha}\, T \left(\vec{u} \cdot \nabla P \right) + Q_L
  \label{eq:energy}
\end{equation}

where \(\sigma^{\prime}\) is the deviatoric stress tensor, \(\rho\) is
density, \(g\) is gravitational acceleration, \(t\) is time,
\(\bar{C}_p\), \(\bar{k}\), \(\bar{\alpha}\) are the reference specific
heat capacity, thermal conductivity, and thermal expansivity,
respectively (see Section \ref{sec:adiabatic-reference-conditions}), and
\(Q_L\) is the latent heat released or absorbed during phase
transitions. Equations \ref{eq:navier-stokes-no-inertia} and
\ref{eq:continuity-compressible} together describe the buoyancy-driven
flow of an isotropic fluid with negligible inertia and Equation
\ref{eq:energy} describes the conduction, advection, and production (or
consumption) of thermal energy \citep{schubert2001}. Note that the
pressure \(P\) in this context is equal to the mean normal stress and is
positive under compression:
\(P = - \frac{\sigma_{xx} + \sigma_{yy} + \sigma_{zz}}{3}\).

The compressible form of the continuity equation (Equation
\ref{eq:continuity-compressible}) is essential for capturing the full
coupling between density changes, pressure and temperature (PT)
variations, and kinetically-controlled phase transitions. This is in
contrast to simplified formulations such as the Boussinesq approximation
or anelastic liquid approximation, which either neglect the time
derivative of density (\(\frac{\partial \rho}{\partial t}\)) entirely or
impose restrictions on where density variations can appear in the
governing equations \citep{gassmoller2020}. By retaining
\(\frac{\partial \rho}{\partial t}\), the compressible continuity
equation enables density changes from kinetically-controlled phase
transitions to influence the flow field anywhere within the model
domain---not just through equilibrium thermodynamics, but through
time-dependent reaction progress. This bidirectional coupling between
flow dynamics and reaction kinetics is critical for accurately
simulating systems where phase transitions occur over finite timescales
comparable to advective timescales.

To solve Equation \ref{eq:continuity-compressible} numerically we adopt
the \emph{projected density approximation}
\citep[PDA,][]{gassmoller2020}, which reformulates the continuity
equation by applying the product rule to
\(\nabla \cdot (\rho\, \vec{u})\) and multiplying both sides by
\(\frac{1}{\rho}\):

\begin{equation}
  \frac{1}{\rho} \frac{\partial \rho}{\partial t} + \nabla \cdot \vec{u} + \left(\frac{1}{\rho} \nabla \rho \right) \cdot \vec{u} = 0
  \label{eq:continuity-expanded}
\end{equation}

The projected density field \(\rho(T, P, X)\) varies with temperature,
pressure, and reaction progress, ensuring that local density changes
arising from both PT variations and phase transitions influence the flow
through buoyancy as well as volumetric expansion and compression. The
phase transitions themselves are modeled using a separate kinetic rate
law (described in Section \ref{sec:reaction-kinetics}), which determines
the reaction progress \(X\) based on local thermodynamic and kinetic
conditions. This makes the PDA ideally suited for our numerical
experiments, which incorporate density changes due to the olivine
\(\Leftrightarrow\) wadsleyite phase transition.

\subsection{Numerical Setup}\label{sec:numerical-setup}

\subsubsection{Adiabatic Reference
Conditions}\label{sec:adiabatic-reference-conditions}

To ensure numerical convergence, we initialized our ASPECT simulations
with reasonable estimates of the pressure-temperature (PT) fields and
material properties in Earth's upper mantle. We began by evaluating
entropy changes over a PT range of 1573--1973 K and 0.001--25 GPa
(Figure \ref{fig:isentrope}) using the Gibbs free energy minimization
software Perple\_X \citep[v.7.0.9,][]{connolly2009}. We assumed a dry
pyrolitic bulk composition after \citet{green1979} and phase equilibria
were evaluated in the Na\(_2\)O-CaO-FeO-MgO-Al\(_2\)O\(_3\)-SiO\(_2\)
(NCFMAS) chemical system with thermodynamic data and solution models of
\citet{stixrude2022}. Equations of state were included for solid
solution phases: olivine, plagioclase, spinel, clinopyroxene,
wadsleyite, ringwoodite, perovskite, ferropericlase, high‐pressure C2/c
pyroxene, orthopyroxene, akimotoite, post‐perovskite, Ca‐ferrite,
garnet, and Na-Al phase.

\begin{figure}
\centering
\includegraphics[width=0.65\linewidth,height=\textheight,keepaspectratio,alt={Entropy (left) and density (right) changes in Earth's upper mantle under thermodynamic equilibrium and hydrostatic stress conditions. Material properties were computed with Perple\_X using the equations of state and thermodynamic data of {[}@stixrude2022{]}. The black box indicates the approximate PT range of our ASPECT simulations, while the white line indicates the isentropic adiabat used to calculate reference material properties.}]{../figs/PYR-material-table.png}
\caption{Entropy (left) and density (right) changes in Earth's upper
mantle under thermodynamic equilibrium and hydrostatic stress
conditions. Material properties were computed with Perple\_X using the
equations of state and thermodynamic data of \citet{stixrude2022}. The
black box indicates the approximate PT range of our ASPECT simulations,
while the white line indicates the isentropic adiabat used to calculate
reference material properties.}\label{fig:isentrope}
\end{figure}

We then determined the mantle adiabat by applying the Newton-Raphson
algorithm to find temperatures corresponding to each pressure that
maintain constant entropy (white line in Figure \ref{fig:isentrope}).
Material properties were evaluated at each PT point along the isentrope
to construct the adiabatic reference conditions shown in Figure
\ref{fig:material-property-profile}. These reference conditions serve
three main purposes: 1) initializing the PT fields and material
properties in our ASPECT simulations (see Section
\ref{sec:initialization-and-boundary-conditions}), 2) updating the
material model during the simulations (see Section
\ref{sec:material-model}), and 3) serving as a basis for computing
``dynamic'' quantities, such as the dynamic temperature
\(\hat{T} = T - \bar{T}\), dynamic pressure \(\hat{P} = P - \bar{P}\),
and dynamic density \(\hat{\rho} = \rho - \bar{\rho}\), that quantify
how much the approximate numerical solution deviates from the adiabatic
reference conditions.

\begin{figure}
\centering
\includegraphics[width=1\linewidth,height=\textheight,keepaspectratio,alt={Reference material properties used in our ASPECT simulations. Profiles were computed using the BurnMan software {[}@cottaar2014; @myhill2023{]} and were based on the equations of state and thermodynamic data of {[}@stixrude2022{]} for pure Mg olivine (ol) and wadsleyite (wd).}]{../figs/material-property-profile.png}
\caption{Reference material properties used in our ASPECT simulations.
Profiles were computed using the BurnMan software
\citep{cottaar2014, myhill2023} and were based on the equations of state
and thermodynamic data of \citet{stixrude2022} for pure Mg olivine (ol)
and wadsleyite (wd).}\label{fig:material-property-profile}
\end{figure}

\subsubsection{Initialization and Boundary
Conditions}\label{sec:initialization-and-boundary-conditions}

We setup our ASPECT simulations within a 900 \(\times\) 600 km
rectangular model domain, initialized with pure Mg olivine and
wadsleyite (Figure \ref{fig:initial-setup}). ``Surface'' PT conditions
of 10 GPa and 1706 K were applied at the top boundary such that the
olivine \(\Leftrightarrow\) wadsleyite transition occurs at
approximately 130--140 km from the top boundary. The initial adiabatic
PT profiles were computed by numerically integrating the following
equations:

\begin{equation}
  \frac{d\bar{T}}{dy} = \frac{\bar{\alpha}\, \bar{T}\, g}{\bar{C}_p}
  \label{eq:adiabatic-temperature}
\end{equation}

\begin{equation}
  \frac{d\bar{P}}{dy} = \bar{\rho}\, g
  \label{eq:adiabatic-pressure}
\end{equation}

where the material properties \(\bar{\rho}\), \(\bar{\alpha}\), and
\(\bar{C}_p\) were determined from the adiabatic reference conditions
shown in Figure \ref{fig:material-property-profile}.

\begin{figure}
\centering
\includegraphics[width=0.45\linewidth,height=\textheight,keepaspectratio,alt={Initial setup for slab (top) and plume (bottom) simulations. The top boundary has a constant ``surface'' PT of 10 GPa and 1706 K such that the olivine \textbackslash Leftrightarrow wadsleyite phase transition (dashed line) occurs at 130--140 km from the top boundary. Thermal anomalies with Gaussian profiles were superimposed on top of the initial adiabatic temperature profile and remained fixed at the inflow boundary. Constant inflow velocities of 5 cm/yr were prescribed parallel to the thermal anomalies. Normal tractions equal to the initial lithostatic pressure profile were enforced on the left, right, and open boundaries to ensure that outflows are driven only by dynamic pressures.}]{../figs/initial-setup.png}
\caption{Initial setup for slab (top) and plume (bottom) simulations.
The top boundary has a constant ``surface'' PT of 10 GPa and 1706 K such
that the olivine \(\Leftrightarrow\) wadsleyite phase transition (dashed
line) occurs at 130--140 km from the top boundary. Thermal anomalies
with Gaussian profiles were superimposed on top of the initial adiabatic
temperature profile and remained fixed at the inflow boundary. Constant
inflow velocities of 5 cm/yr were prescribed parallel to the thermal
anomalies. Normal tractions equal to the initial lithostatic pressure
profile were enforced on the left, right, and open boundaries to ensure
that outflows are driven only by dynamic
pressures.}\label{fig:initial-setup}
\end{figure}

For the initial temperature field, thermal anomalies with amplitudes of
\(\pm\) 500 K were superimposed on the adiabatic temperature profile.
These anomalies were defined as smooth linear features with Gaussian
cross-sections (15 km half-width) and tanh tapered ends (5 km taper
length) to avoid sharp discontinuities. Slabs extended 100 km
horizontally and 100 km downward from the top boundary; plumes extended
450 km upward from the bottom boundary. Velocity boundary conditions
prescribed constant inflow of 5 cm/yr parallel to the thermal anomalies,
tapering smoothly to zero at the thermal anomaly edges with the same
Gaussian profile. Zero horizontal velocities were imposed at the side
boundaries (\(\vec{u}_x\) = 0). The full functional form of the
Gaussian-tanh thermal anomalies and velocity boundary conditions are
available within the accompanying data repository (see Data Availability
statement for details).

Stress boundary conditions on the left and right boundaries enforced a
normal traction equal to the lithostatic pressure profile computed in
Equation \ref{eq:adiabatic-pressure}. No tangential (shear) stresses
were applied to the side boundaries, so they approximated impermeable,
free-slip surfaces under hydrostatic confinement. Open boundaries
(bottom for slabs, top for plumes) were assigned a constant normal
traction equal to the initial lithostatic pressure at the respective
boundary \(\sigma_{yy}\) = \(\bar{P}(y)\). These stress conditions allow
outflow to occur freely at the top or bottom boundaries (for plumes
versus slabs), driven only by dynamic pressure variations associated
with convection and/or volume changes during the olivine
\(\Leftrightarrow\) wadsleyite phase transition.

\subsubsection{Material Model}\label{sec:material-model}

\paragraph{Material Properties}\label{sec:material-properties}

Material properties were updated during our ASPECT simulations by
referencing the adiabatic reference conditions shown in Figure
\ref{fig:material-property-profile}. Except for density, material
properties received no PT corrections, effectively assuming that
deviations from the adiabatic reference conditions were negligible. For
density, however, we applied a dynamic PT correction through a
first-order Taylor expansion \citep{jarvis1980, gassmoller2020}:

\begin{equation}
  \rho \approx \bar{\rho} + \left(\frac{\partial \bar{\rho}}{\partial P} \right)_T \Delta P + \left(\frac{\partial \bar{\rho}}{\partial T} \right)_P \Delta T
  \label{eq:density-ala-expansion}
\end{equation}

Equation \ref{eq:density-ala-expansion} is rewritten using standard
thermodynamic relations
\(\beta = \frac{1}{\rho} \left(\frac{\partial \rho}{\partial P}\right)_T\)
and
\(\alpha = -\frac{1}{\rho} \left(\frac{\partial \rho}{\partial T}\right)_P\)
to obtain the expression:

\begin{equation}
  \rho \approx \bar{\rho} \left(1 + \bar{\beta}\, \hat{P} - \bar{\alpha}\, \hat{T} \right)
  \label{eq:density-ala}
\end{equation}

where \(\bar{\rho}\), \(\, \bar{\beta}\), \(\, \bar{\alpha}\), are the
adiabatic reference density, compressibility, and thermal expansivity,
respectively, and \(\Delta P = \hat{P} = P - \bar{P}\) and
\(\Delta T = \hat{T} = T - \bar{T}\) are the dynamic PT. Note that the
reference thermal conductivity \(\bar{k}\) = 4.0 W m\(^{-1}\)K\(^{-1}\)
is constant in all our numerical experiments.

\paragraph{Reaction Kinetics}\label{sec:reaction-kinetics}

The kinetics of the olivine \(\Leftrightarrow\) wadsleyite phase
transition were governed entirely by interface-controlled growth, as
nucleation was assumed to saturate rapidly and did not limit the
reaction \citep{cahn1956}. Following \citet{cahn1956}, the transformed
volume fraction is given by:

\begin{equation}
  X = 1 - \exp\!\left(-N\, \dot{x}\, t \right)
  \label{eq:volume-fraction}
\end{equation}

where \(X\) is the volume fraction of the product phase (olivine or
wadsleyite), \(N\) is a geometric factor that accounts for nucleation
sites, \(\dot{x}\) is the growth rate, and \(t\) is the elapsed time
after site saturation. For inter-crystalline grain-boundary controlled
growth, \(N = 6.67/d\), where \(d\) is grain size.

Since we assumed interface-controlled growth kinetics, the following
expression determined the overall reaction rate \citep{hosoya2005}:

\begin{equation}
  \dot{x} = \Gamma\, T\, C_\mathrm{OH}^n\, \exp\!\left(-\frac{H^{\ast} + P V^{\ast}}{R\, T}\right) \left(1 - \exp\!\left[-\frac{\Delta G}{R\, T}\right] \right)
  \label{eq:growth-rate}
\end{equation}

where \(\Gamma\) is the growth rate prefactor, \(C_\mathrm{OH}\) is the
concentration of water in the reactant phase, \(n\) is the water content
exponent, \(H^{\ast}\) is activation enthalpy, \(V^{\ast}\) is
activation volume, \(P\) is pressure, \(T\) is temperature, \(R\) is the
gas constant, and \(\Delta G\) is the Gibbs free energy difference
between olivine and wadsleyite, which is approximated by:

\begin{equation}
  \Delta G \approx \Delta \bar{G} + \hat{P}\, \Delta \bar{V} - \hat{T}\, \Delta \bar{S}
  \label{eq:gibbs}
\end{equation}

where \(\Delta \bar{G}\), \(\Delta \bar{V}\), and \(\Delta \bar{S}\) are
the molar Gibbs free energy, volume, and entropy differences between
olivine and wadsleyite along the adiabatic reference profile (Figure
\ref{fig:thermodynamic-property-profile}), respectively, and \(\hat{P}\)
and \(\hat{T}\) are the dynamic PT.

This formulation represents a significant thermodynamic simplification.
We define Gibbs free energy based on mean normal stress (pressure), even
though our simulations involve non-hydrostatic stress conditions where
deviatoric stresses and normal stresses at individual grain interfaces
directly influence reaction pathways and chemical potentials
\citep{wheeler2015, wheeler2020}. Mean stress can provide a reasonable
approximation under certain conditions \citep{wheeler2015b}. A rigorous
treatment will require consideration of stress-dependent chemical
potentials that account for the full stress tensor's influence on phase
stability and reaction kinetics (work in progress).

Setting aside these important microstructural effects for now, the time
evolution of the olivine \(\Leftrightarrow\) wadsleyite phase transition
is fully described by the interplay of pressure, temperature, and
kinetic parameters applied to the interface-controlled growth model,
without explicit consideration of nucleation kinetics
\citep{hosoya2005, faccenda2017}. The macro-scale olivine
\(\Leftrightarrow\) wadsleyite reaction rate was therefore computed by
taking the time derivative of Equation \ref{eq:volume-fraction}:

\begin{equation}
  \frac{dX}{dt} = \dot{X} = N\, \dot{x}\, \left(1 - X \right)
  \label{eq:reaction-rate-short}
\end{equation}

Since all of the kinetic parameters \(N\), \(\Gamma\), and
\(C_\mathrm{OH}^n\) ultimately scale the reaction rate (Equations
\ref{eq:growth-rate}--\ref{eq:reaction-rate-short}), varying them
independently adds little scientific value. Instead, we simplified our
numerical implementation of Equation \ref{eq:reaction-rate-short} by
combining the parameters \(N\), \(\Gamma\), and \(C_\mathrm{OH}^n\) into
a single kinetic prefactor
\(Z = \frac{6.67}{d}\, \Gamma\, C_\mathrm{OH}^n\). Thus, the full
expression for the reaction rate became:

\begin{equation}
  \dot{X} = Z\, T\, \exp\!\left(-\frac{H^{\ast} + P V^{\ast}}{R\, T}\right) \left(1 - \exp\!\left[-\frac{\Delta G}{R\, T}\right] \right)\, \left(1 - X \right)
  \label{eq:reaction-rate}
\end{equation}

The range of kinetic prefactors \(Z\) used in our numerical experiments
(3.0e0--7.0e7 K s\(^{-1}\)) was determined by holding \(\Gamma\) =
\(\exp\!\left(-18\right)\) m s\(^{-1}\) K\(^{-1}\)
ppm\(_\mathrm{OH}^{-n}\), \(H^{\ast}\) = 274 kJ mol\(^{-1}\),
\(V^{\ast}\) = 3.0e-6 m\(^3\) mol\(^{-1}\), and \(n\) = 3.2 constant,
while varying water content \(C_\mathrm{OH}\) from 50--5000 ppm and
grain size \(d\) from 1--10 mm. These water contents and grain sizes are
consistent with the experimental conditions of \citet{hosoya2005},
previous numerical studies of metastable olivine wedges
\citep{rubie1994}, and typical grain sizes of upper mantle xenoliths
\citep[\textasciitilde3--10 mm,][]{karato1984, karato2008}. Thus, our
experiments approximate kinetic conditions ranging from slow kinetics in
dry rocks with large grain sizes (50 ppm OH; 10 mm; \(Z\) = 3.0e0 K
s\(^{-1}\)) to fast kinetics in hydrated rocks with small grain sizes
(5000 ppm OH; 1 mm; \(Z\) = 7.0e7 K s\(^{-1}\)).

\begin{figure}
\centering
\includegraphics[width=0.8\linewidth,height=\textheight,keepaspectratio,alt={Reference thermodynamic properties used in our ASPECT simulations. Profiles were computed using the same methods as described in Figure  (see Section ).}]{../figs/thermodynamic-property-profile.png}
\caption{Reference thermodynamic properties used in our ASPECT
simulations. Profiles were computed using the same methods as described
in Figure \ref{fig:material-property-profile} (see Section
\ref{sec:adiabatic-reference-conditions}).}\label{fig:thermodynamic-property-profile}
\end{figure}

\paragraph{Operator Splitting}\label{sec:operator-splitting}

Since the reaction rate \(\dot{X}\) was faster than the advection
timescale in our ASPECT simulations, we employed a first-order operator
splitting scheme to decouple advection from interface-controlled olivine
\(\Leftrightarrow\) wadsleyite growth kinetics. In this approach, the
transformed volume fraction \(X\) was updated in two sequential steps
within each overall time step \(\Delta t\):

\begin{equation}
  \frac{\partial X}{\partial t} + \vec{u} \cdot \nabla X = 0
  \label{eq:composition}
\end{equation}

\begin{enumerate}
\def\labelenumi{\arabic{enumi}.}
\tightlist
\item
  \textbf{Reaction step:} Starting from \(X^{t}\), the reaction rate
  equation (Equation \ref{eq:reaction-rate}) was integrated forward in
  time over the full advection timestep \(\Delta t\), yielding an
  intermediate composition \(X^{\ast}\). This integration was performed
  using a sequence of internally determined reaction substeps
  \(\delta t \le \Delta t\), such that the reaction operator was applied
  repeatedly until the cumulative integration time reached \(\Delta t\).
\item
  \textbf{Advection step:} Starting from \(X^{\ast}\), the material
  transport term \(\left(\vec{u} \cdot \nabla X \right)\) in Equation
  \ref{eq:composition} was solved over the advection time interval
  \(\Delta t\), neglecting phase changes, to obtain the updated
  composition \(X^{t+1}\).
\end{enumerate}

In our simulations, the reaction step was solved using the ARKODE
package of the SUNDIALS library \citep{reynolds2023, hindmarsh2005},
which employs an adaptive-step additive Runge-Kutta method
\citep{aspectmanual}. The reaction substep size \(\delta t\) was
determined dynamically by ARKODE to satisfy a specified relative
tolerance during the reaction step (set to \(10^{-6}\) in this study).
As a result, the reaction operator was subcycled as needed within each
global timestep \(\Delta t\), ensuring that fast solid-state reaction
kinetics were accurately captured without imposing overly restrictive
global timesteps.

\subsubsection{Rheological Model}\label{sec:rheological-model}

We use a temperature-dependent viscosity following an Arrhenius law:

\begin{equation}
  \eta = A \exp\!\left(\frac{E^{\ast}}{R\,T}\right)
  \label{eq:arrhenius-viscosity}
\end{equation}

with viscosity prefactor \(A\), activation energy \(E^{\ast}\), gas
constant \(R\), and absolute temperature \(T\), which is decomposed into
an adiabatic reference temperature \(\bar{T}\) and a perturbation
\(\hat{T}\), \(T=\bar{T}+\hat{T}\). Linearizing the Arrhenius relation
through a first-order Taylor expansion about \(\bar{T}\) yields:

\begin{equation}
  \eta \approx A \exp\!\left(\frac{E^{\ast}}{R\,\bar{T}}\right) \exp\!\left(-\frac{E^{\ast}}{R\,\bar{T}}\frac{\hat{T}}{\bar{T}}\right)
  \label{eq:arrhenius-viscosity-expanded}
\end{equation}

By defining a reference background viscosity as:

\begin{equation}
  \bar{\eta} = A \exp\!\left(\frac{E^{\ast}}{R\,\bar{T}}\right)
  \label{eq:background-viscosity}
\end{equation}

we arrive at a rheological model where viscosity varies exponentially
with thermal perturbations about an adiabatic reference profile:

\begin{equation}
  \eta \approx \bar{\eta}\,\exp\!\left(-B\,\frac{\hat{T}}{\bar{T}}\right) \qquad B=\frac{E^{\ast}}{R\,\bar{T}}
  \label{eq:rheological-model}
\end{equation}

In our simulations, we prescribed a uniform background viscosity
\(\bar{\eta}\) = 10\(^{21}\) Pa s throughout the upper mantle
\citep[e.g.,][]{karato2008, ranalli1995}, implicitly assuming that
\(E^{\ast}\) was temperature-dependent and varied slightly with depth
(Equation \ref{eq:background-viscosity}). We varied the rheological
activation factor \(B\) between 2 (low thermal sensitivity) and 10 (high
thermal sensitivity). For our prescribed thermal anomalies of
\(\hat{T}\) = \(\pm\) 500 K, these \(B\) values correspond to
approximately 2--20\(\times\) stronger slabs or 2--20\(\times\) weaker
plumes relative to the background mantle for \(B\) = 2 and 10,
respectively. Thus, our numerical experiments explore a range of
viscosity contrasts between thermal anomalies and background adiabatic
reference conditions, with the exponential rheological model creating
symmetric viscosity variations for heating and cooling.

\subsubsection{Numerical Stabilization of Dynamic Pressure
Oscillations}\label{sec:numerical-stability}

A significant numerical limitation emerges from the coupled feedback
between our kinetic rate law (Equations
\ref{eq:volume-fraction}--\ref{eq:reaction-rate}) and the buildup of
dynamic pressure in the fully compressible continuity equation (Equation
\ref{eq:continuity-compressible}). At sufficiently high rheological
contrasts (\(B\) \(\gtrsim\) 4), cold slabs develop internal dynamic
pressures that can exceed several hundred MPa. These dynamic pressure
perturbations accelerate the forward reaction (through the Gibbs free
energy term in Equation \ref{eq:reaction-rate}), causing rapid local
density changes. These density changes then alter the pressure field
through the coupled continuity and momentum equations (Equations
\ref{eq:navier-stokes-no-inertia}--\ref{eq:continuity-compressible}),
which in turn drives the reverse reaction. This positive feedback loop
manifests as spurious pressure waves propagating through the slab
interior.

To address this issue, we adopt an approach similar to
\citet{gassmoller2020} and exclude the dynamic pressure contribution
from the Gibbs free energy calculation (Equation \ref{eq:gibbs}) while
retaining it in the Arrhenius term in Equation \ref{eq:reaction-rate}
and density formulation (Equation \ref{eq:density-ala}). In practice,
this means replacing
\(\Delta G \approx \Delta \bar{G} + \hat{P}\, \Delta \bar{V} - \hat{T}\, \Delta \bar{S}\)
with \(\Delta G \approx \Delta \bar{G} - \hat{T}\, \Delta \bar{S}\) in
our kinetic rate law. As discussed by \citet{gassmoller2020}, this
approximation is justified because density variations from dynamic
pressure are typically small compared to those from compositional
heterogeneities and temperature variations in mantle convection.
However, this treatment does introduce a limitation. In scenarios where
dynamic pressure effects dominate over thermal and compositional
contributions, such as in regions with extreme viscosity contrasts or
near phase boundaries with large \(\Delta \bar{V}\), our kinetic model
may underestimate the thermodynamic driving force for phase transitions.
Nevertheless, this compromise ensures numerical stability while
preserving the essential physics of kinetically controlled phase
transitions coupled to compressible flow.

\section{Results}\label{sec:results}

\subsection{Simulation Snapshots: Slabs and
Plumes}\label{sec:simulation-snapshots}

Figures \ref{fig:slab-composition-set2} and
\ref{fig:plume-composition-set2} illustrate how reaction rates impact
dynamic flow patterns and shape the 410 discontinuity. These snapshots,
taken after 100 Ma of evolution, provide visual context for the
quantitative analysis in Section \ref{sec:410-displacement-width}.

In slab simulations, ultra-sluggish kinetics (Figure
\ref{fig:slab-composition-set2}: top row) allow metastable olivine to
persist deep into the transition zone. This inhibition causes the slab
to stagnate and pond, depressing the 410. Within the cold, metastable
olivine region, Gibbs free energy accumulates and wadsleyite saturation
remains low until the thermodynamic driving force overcomes kinetic
barriers. Once this threshold is reached, the olivine
\(\Leftrightarrow\) wadsleyite reaction rapidly completes, producing a
sharp 410 that is displaced downwards by tens of kilometers.

At intermediate reaction rates (Figure \ref{fig:slab-composition-set2};
middle row), the olivine \(\Leftrightarrow\) wadsleyite reaction still
lags but is fast enough to limit widespread olivine metastability and
avoid total slab stagnation. The resulting 410 is broad and diffuse, as
density and seismic velocity contrasts gradually fade with depth. This
moderately-sluggish kinetic regime produces complex 410 structures
through intermediate reaction rates, incomplete slab stagnation, and
deflected flow patterns. However, when reaction rates are sufficiently
fast to maintain quasi-equilibrium conditions (Figure
\ref{fig:slab-composition-set2}; bottom row), the 410 sharpens and
shifts upwards as expected from equilibrium thermodynamics. Under this
fast kinetic regime, rapid wadsleyite growth within the slab allows
continuous slab descent through the 410 without hesitation.

In plume simulations, thermal effects dominate mantle flow dynamics and
410 structure. Even under ultra-sluggish kinetics (Figure
\ref{fig:plume-composition-set2}; top row), the high temperatures of
upwellings prevent significant wadsleyite metastability. The olivine
\(\Leftrightarrow\) wadsleyite transition proceeds rapidly, maintaining
thin, sharp 410 interfaces and strong density and seismic contrasts.
Although ultra-sluggish kinetics slightly broaden and uplift the 410,
reducing buoyancy contrasts, plume structures remain vertically coherent
across the full range of tested kinetic prefactors \(Z\) (Figure
\ref{fig:plume-composition-set2}).

Altogether, these simulations demonstrate that in cold environments,
kinetics strongly influence slab dynamics and control whether the 410
appears as a diffuse, low-amplitude feature or as a sharp, high-contrast
seismic boundary. In contrast, thermal effects dominate in hot plume
environments, producing stable, sharply defined 410s that are largely
independent of tested kinetic prefactors.

\begin{figure}
\centering
\includegraphics[width=1\linewidth,height=\textheight,keepaspectratio,alt={Slab simulations with intermediate strength contrasts (B = 4, viscosity contrast \textasciitilde3\textbackslash times) showing ultra-sluggish (top row: Z = 3.0e0 K s\^{}\{-1\}), intermediate (middle row: Z = 4.7e2 K s\^{}\{-1\}), and quasi-equilibrium (bottom row: Z = 7.0e7 K s\^{}\{-1\}) kinetic regimes after 100 Ma evolution. Panels show dynamic temperature \textbackslash hat\{T\} (left column), dynamic density \textbackslash hat\{\textbackslash rho\} (middle column), and pressure-wave velocity V\_p (right column). Thin lines highlight the 10\% and 90\% wadsleyite volume fraction contours (X = 0.1 and 0.9). The 410 displacement is defined as the difference between the depth at X = 0.9 and the nominal equilibrium olivine \textbackslash Leftrightarrow wadsleyite transition depth, while the 410 width is defined as the difference between depths at X = 0.9 and X = 0.1 (see Supplementary Information for details).}]{../figs/simulation/compositions/slab-Z3.0e0-B4-Z4.7e2-B4-Z7.0e7-B4-set2-composition-0010.png}
\caption{Slab simulations with intermediate strength contrasts (\(B\) =
4, viscosity contrast \textasciitilde3\(\times\)) showing ultra-sluggish
(top row: \(Z\) = 3.0e0 K s\(^{-1}\)), intermediate (middle row: \(Z\) =
4.7e2 K s\(^{-1}\)), and quasi-equilibrium (bottom row: \(Z\) = 7.0e7 K
s\(^{-1}\)) kinetic regimes after 100 Ma evolution. Panels show dynamic
temperature \(\hat{T}\) (left column), dynamic density \(\hat{\rho}\)
(middle column), and pressure-wave velocity \(V_p\) (right column). Thin
lines highlight the 10\% and 90\% wadsleyite volume fraction contours
(\(X\) = 0.1 and 0.9). The 410 displacement is defined as the difference
between the depth at X = 0.9 and the nominal equilibrium olivine
\(\Leftrightarrow\) wadsleyite transition depth, while the 410 width is
defined as the difference between depths at X = 0.9 and X = 0.1 (see
Supplementary Information for
details).}\label{fig:slab-composition-set2}
\end{figure}

\begin{figure}
\centering
\includegraphics[width=1\linewidth,height=\textheight,keepaspectratio,alt={Plume simulations with intermediate strength contrasts (B = 4, viscosity contrast \textasciitilde1/3\textbackslash times) showing ultra-sluggish (top row: Z = 3.0e0 K s\^{}\{-1\}), intermediate-sluggish (middle row: Z = 4.7e2 K s\^{}\{-1\}), and quasi-equilibrium (bottom row: Z = 7.0e7 K s\^{}\{-1\}) kinetic regimes after 100 Ma evolution. Panels show dynamic temperature \textbackslash hat\{T\} (left column), dynamic density \textbackslash hat\{\textbackslash rho\} (middle column), and pressure-wave velocity V\_p (right column). Thin lines highlight the 10\% and 90\% wadsleyite volume fraction contours (X = 0.1 and 0.9). The 410 displacement is defined as the difference between the depth at X = 0.9 and the nominal equilibrium olivine \textbackslash Leftrightarrow wadsleyite transition depth, while the 410 width is defined as the difference between depths at X = 0.9 and X = 0.1 (see Supplementary Information for details).}]{../figs/simulation/compositions/plume-Z3.0e0-B4-Z4.7e2-B4-Z7.0e7-B4-set2-composition-0010.png}
\caption{Plume simulations with intermediate strength contrasts (\(B\) =
4, viscosity contrast \textasciitilde1/3\(\times\)) showing
ultra-sluggish (top row: \(Z\) = 3.0e0 K s\(^{-1}\)),
intermediate-sluggish (middle row: \(Z\) = 4.7e2 K s\(^{-1}\)), and
quasi-equilibrium (bottom row: \(Z\) = 7.0e7 K s\(^{-1}\)) kinetic
regimes after 100 Ma evolution. Panels show dynamic temperature
\(\hat{T}\) (left column), dynamic density \(\hat{\rho}\) (middle
column), and pressure-wave velocity \(V_p\) (right column). Thin lines
highlight the 10\% and 90\% wadsleyite volume fraction contours (\(X\) =
0.1 and 0.9). The 410 displacement is defined as the difference between
the depth at X = 0.9 and the nominal equilibrium olivine
\(\Leftrightarrow\) wadsleyite transition depth, while the 410 width is
defined as the difference between depths at X = 0.9 and X = 0.1 (see
Supplementary Information for
details).}\label{fig:plume-composition-set2}
\end{figure}

\subsection{Structure of the 410: Displacement and
Width}\label{sec:410-displacement-width}

Figure \ref{fig:410-structure} summarizes the quantitative relationships
between 410 structure and the maximum reaction rate
\(\dot{X}_{\mathrm{max}}\) evaluated in slab and plume simulations after
100 Ma of evolution. The results reveal fundamentally different
responses of plumes and slabs to reaction kinetics. See the
Supplementary Information for the full set of experimental results and
technical details describing the 410 structural measurements.

\begin{figure}
\centering
\includegraphics[width=0.7\linewidth,height=\textheight,keepaspectratio,alt={Measured 410 displacement and width versus maximum reaction rates \textbackslash dot\{X\}\_\{\textbackslash mathrm\{max\}\} in plume and slab simulations with intermediate strength contrasts (B = 4) after 100 Ma. Structure of the 410 near plumes (left column) shows minimal dependence on \textbackslash dot\{X\}\_\{\textbackslash mathrm\{max\}\}, with both displacement and width remaining nearly constant across seven orders of magnitude variation in \textbackslash dot\{X\}\_\{\textbackslash mathrm\{max\}\}. In contrast, 410 structure near slabs (right column) changes distinctly across three kinetic regimes: (1) quasi-equilibrium at high \textbackslash dot\{X\}\_\{\textbackslash mathrm\{max\}\}, where 410 widths are narrow and elevated; (2) an intermediate regime where decreasing reaction rates \textbackslash dot\{X\}\_\{\textbackslash mathrm\{max\}\} progressively widen and deepen the 410; and (3) an ultra-sluggish regime at low \textbackslash dot\{X\}\_\{\textbackslash mathrm\{max\}\}, where the 410 narrows while deepening, and slabs completely stall and pond.}]{../figs/410-structure-B4.png}
\caption{Measured 410 displacement and width versus maximum reaction
rates \(\dot{X}_{\mathrm{max}}\) in plume and slab simulations with
intermediate strength contrasts (\(B\) = 4) after 100 Ma. Structure of
the 410 near plumes (left column) shows minimal dependence on
\(\dot{X}_{\mathrm{max}}\), with both displacement and width remaining
nearly constant across seven orders of magnitude variation in
\(\dot{X}_{\mathrm{max}}\). In contrast, 410 structure near slabs (right
column) changes distinctly across three kinetic regimes: (1)
quasi-equilibrium at high \(\dot{X}_{\mathrm{max}}\), where 410 widths
are narrow and elevated; (2) an intermediate regime where decreasing
reaction rates \(\dot{X}_{\mathrm{max}}\) progressively widen and deepen
the 410; and (3) an ultra-sluggish regime at low
\(\dot{X}_{\mathrm{max}}\), where the 410 narrows while deepening, and
slabs completely stall and pond.}\label{fig:410-structure}
\end{figure}

In plume simulations, the 410 shows little dependence on
\(\dot{X}_{\mathrm{max}}\). Its structure remains nearly constant across
seven orders of magnitude variation in \(\dot{X}_{\mathrm{max}}\), with
consistent displacements of 24 km and widths between 2--9 km. The only
exception is a few ultra-sluggish kinetic models where displacements
decrease to 18--21 km and widths increase to 12--21 km. The weak
dependence of 410 structure on \(\dot{X}_{\mathrm{max}}\) reflects the
strong thermal control of the reaction front in upwellings, where high
temperatures promote rapid wadsleyite \(\Leftrightarrow\) olivine
transition---maintaining a sharp discontinuity regardless of the kinetic
prefactor \(Z\) applied to the interface-controlled olivine
\(\Leftrightarrow\) wadsleyite growth model (Equation
\ref{eq:reaction-rate}). Similarly, variations in rheological strength
contrast (controlled by the \(B\) parameter in Equation
\ref{eq:rheological-model}) produce negligible changes to plume
morphology or 410 structure, as the thermally-dominated behavior is
largely insensitive to viscosity variations (see Supplementary
Information for examples).

In slab simulations, the 410 exhibits distinct structural changes across
three kinetic regimes. At high reaction rates (\(Z\) \(\gtrsim\) 1.8e5 K
s\(^{-1}\); \(\dot{X}_{\mathrm{max}}\) \(\gtrsim\) 2.2 Ma\(^{-1}\)), the
olivine \(\Leftrightarrow\) wadsleyite transition remains near
thermodynamic equilibrium, producing a narrow 410 (\(\lesssim\) 10 km),
displaced 27--39 km upwards within the slab's inner core. As
\(\dot{X}_{\mathrm{max}}\) decreases (2.0e2 \(\lesssim\) \(Z\)
\(\lesssim\) 1.8e5 K s\(^{-1}\); 0.07 \(\lesssim\)
\(\dot{X}_{\mathrm{max}}\) \(\lesssim\) 2.2 Ma\(^{-1}\)), the 410
deepens and widens, approximating a log-linear relationship where
reductions in \(\dot{X}_{\mathrm{max}}\) progressively broaden the
reaction front. This intermediate kinetic regime corresponds to a
partially inhibited olivine \(\Leftrightarrow\) wadsleyite phase
transition that proceeds slowly, hindering downward flow without
complete slab stagnation.

At the lowest reaction rates (\(Z\) \(\lesssim\) 2.0e2 K s\(^{-1}\);
\(\dot{X}_{\mathrm{max}}\) \(\lesssim\) 0.07 Ma\(^{-1}\)), a third
kinetic regime emerges. While the 410 is displaced downwards, its width
narrows with further reductions in \(\dot{X}_{\mathrm{max}}\). This
ultra-sluggish kinetic regime reflects a transition to strong
disequilibrium conditions and complete slab stagnation. As metastable
olivine ponds and is pushed deeper, increasing pressure drives growing
thermodynamic disequilibrium, eventually triggering rapid transformation
within a narrow depth range while surrounding regions remain kinetically
frozen. The apparent 410 sharpening results from this
pressure-controlled reaction front; only material pushed sufficiently
deep accumulates enough thermodynamic driving force to react, producing
a sharp seismic discontinuity where this critical pressure is reached.

Rheological strength contrasts substantially modulate slab dynamics and
410 structure through interactions with kinetic effects (Figure
\ref{fig:410-structure-comp}). Higher \(B\) values in Equation
\ref{eq:rheological-model} increase the viscosity contrast between the
cold slab and surrounding mantle, producing a more coherent slab that
penetrates the 410 with less internal deformation. Stronger slabs sink
more sub-horizontally through the 410, limiting their vertical descent
rates. Therefore, as \(B\) increases, the 410 sharpens because material
traverses the phase transition zone more slowly, allowing greater time
for reaction progress despite sluggish kinetics. In contrast, lower
\(B\) values permit greater internal deformation and faster vertical
descent, which amplifies kinetic effects and produces broader phase
transition zones (see Supplementary Information for examples).

\begin{figure}
\centering
\includegraphics[width=0.7\linewidth,height=\textheight,keepaspectratio,alt={Variation in 410 structure and slab descent rate across kinetic and rheological regimes. Panels show maximum reaction rate (top left), maximum vertical velocity (top right), 410 width (bottom left), and 410 displacement (bottom right) as functions of the kinetic prefactor Z (horizontal axis, log scale) and rheological activation factor B (vertical axis). Each colored tile represents the measured value within a simulation after 100 Ma, and black/white lines delineate transitions between regime behaviors. The precise location of these transitions depends on B, where increasing rheological contrast progressively shifts the regime boundaries towards lower Z (more sluggish kinetic conditions). Text labels represent the observed kinetic regimes: (1) quasi-equilibrium, (2) intermediate, and (3) ultra-sluggish. The complete dataset is given in the Supplementary Information.}]{../figs/410-structure-comp.png}
\caption{Variation in 410 structure and slab descent rate across kinetic
and rheological regimes. Panels show maximum reaction rate (top left),
maximum vertical velocity (top right), 410 width (bottom left), and 410
displacement (bottom right) as functions of the kinetic prefactor \(Z\)
(horizontal axis, log scale) and rheological activation factor \(B\)
(vertical axis). Each colored tile represents the measured value within
a simulation after 100 Ma, and black/white lines delineate transitions
between regime behaviors. The precise location of these transitions
depends on \(B\), where increasing rheological contrast progressively
shifts the regime boundaries towards lower \(Z\) (more sluggish kinetic
conditions). Text labels represent the observed kinetic regimes: (1)
quasi-equilibrium, (2) intermediate, and (3) ultra-sluggish. The
complete dataset is given in the Supplementary
Information.}\label{fig:410-structure-comp}
\end{figure}

In summary, 410 structure near plumes is regulated by thermal effects
near thermodynamic equilibrium, whereas 410 structure near slabs
exhibits distinct kinetic thresholds and non-linear scaling between its
width, displacement, and the reaction rate \(\dot{X}\). These
contrasting behaviors underscore the differing roles of kinetics in hot
versus cold mantle environments and imply that the 410 beneath slabs can
transition abruptly between thermodynamically- and
kinetically-controlled regimes as reaction rates decrease.

\section{Discussion}\label{sec:discussion}

Our numerical simulations show that reaction kinetics exert a
first-order control on the structure and seismic expression of the 410
discontinuity. The results presented in Section \ref{sec:results}
demonstrate that plume and slab dynamics respond in systematically
different ways. Plumes are insensitive to kinetics due to high
temperatures, whereas slabs show three distinct kinetic regimes with
thresholded behavior (Figure \ref{fig:410-structure}). Rheological
strength contrasts, controlled by the activation factor \(B\), further
modulate these kinetic effects by altering slab geometry and descent
rate (Figure \ref{fig:410-structure-comp}). The implications of such
contrasting relationships in slabs versus plumes are discussed below.

\subsection{Uncertainties and Model
Limitations}\label{sec:uncertainties-and-model-limitations}

The primary quantitative uncertainty in our analysis stems from the
kinetic prefactor \(Z\) in the interface-controlled olivine
\(\Leftrightarrow\) wadsleyite growth model (Equation
\ref{eq:reaction-rate}), which spans several orders of magnitude
reflecting variable water contents (50--5000 ppm) and grain sizes (1--10
mm). Laboratory studies also reveal large uncertainties in kinetic
parameters \(n\), \(\Gamma\), \(H^{\ast}\), and \(V^{\ast}\) that depend
strongly on water content, grain size, Mg-Fe composition, and
microstructural evolution
\citep{rubie1994, liu1998, kubo2004, hosoya2005, perrillat2013, ledoux2023}.
Our simulations therefore explore only a limited subset of potential
reaction rates in Earth's upper mantle.

Our kinetic model also assumes instantaneous nucleation site saturation
followed by interface-controlled growth (Equations
\ref{eq:volume-fraction}--\ref{eq:reaction-rate}), thereby neglecting
nucleation kinetics. While often justified because nucleation rates
occur too rapidly for reliable measurement
\citep{hosoya2005, kubo2004, faccenda2017, perrillat2016}, recent
in-situ X-ray and acoustic studies \citep{ohuchi2022, ledoux2023}
document complex nucleation-growth microstructures that can limit net
reaction rates under some PT conditions. By assuming saturated
nucleation, we generally overestimate reaction rates and consequently
underestimate the degree of olivine metastability and its influence on
flow dynamics and 410 structure.

Beyond these kinetic uncertainties, our simulations incorporate several
compositional and rheological simplifications. Assuming pure Mg
endmembers neglects Fe-partitioning and minor element effects that can
shift equilibrium depths by \textasciitilde10--20 km and alter kinetics
\citep{katsura2004, perrillat2013, perrillat2016}. We also neglect
non-hydrostatic stress effects on microstructures and solid-state
reactions \citep{wheeler2014, wheeler2018, wheeler2020}. Our simple
temperature-dependent viscosity (Equations
\ref{eq:arrhenius-viscosity}--\ref{eq:rheological-model}) omits several
strain-dependent processes that could affect both flow dynamics and
reaction kinetics. Notably, accumulated strain could influence 410
topography through multiple pathways: grain size reduction would
accelerate kinetics by increasing nucleation site density (affecting
\(N\) in Equation \ref{eq:reaction-rate-short}), strain localization
could similarly create zones of enhanced reaction rates, and
strain-weakening could modify slab descent trajectory and residence time
in the transition zone. While incorporating strain-weakening rheologies,
grain-size evolution, and plastic deformation would provide a more
complete treatment \citep{karato2001}, such additions would
substantially expand an already large parameter space and introduce
additional uncertainties in constitutive relationships.

Finally, our 2D simulations neglect 3D slab geometry and dynamics. While
we expect the same kinetic regimes to occur in 3D, the transitions
between regimes will likely shift as slabs behave differently in 3D
versus 2D \citep{sime2024}. Despite these combined limitations in
kinetic parameters, compositional complexity, rheological formulation,
and dimensionality, our results capture first-order effects governing
410 structure.

\subsection{Implications for Subduction
Dynamics}\label{sec:implications-for-subduction-dynamics}

The three kinetic regimes identified in our slab
simulations---quasi-equilibrium, intermediate, and
ultra-sluggish---emerge from feedbacks between kinetically controlled
phase transitions and slab strength (Figure
\ref{fig:410-structure-comp}). However, seismic tomography shows little
to no evidence for widespread slab ponding at the 410 \citep{fukao2013},
providing an observational constraint on plausible kinetic behavior in
natural systems. In particular, the ultra-sluggish regime (\(\dot{X}\)
\(\lesssim\) 0.07 Ma\(^{-1}\)), which leads to complete slab stagnation
at the 410 in our simulations, is inconsistent with global seismic
observations of continuous slab descent through the 410. This implies
that most subduction zones experience sufficiently rapid olivine
\(\Leftrightarrow\) wadsleyite transformation to avoid long-lived
stagnation and ponding at the 410.

Meanwhile, the occurrence of deep earthquakes often attributed to
transformational faulting could indicate metastable olivine persistence
in various subduction zones
\citep{green1995, kirby1996, ishii2021, ohuchi2022, sindhusuta2025} and
thus be compatible with the intermediate kinetic regime in our
simulations (0.07 \(\lesssim\) \(\dot{X}\) \(\lesssim\) 2.2
Ma\(^{-1}\)). This regime would generate localized buoyant regions that
resist but do not prevent downward flow. Assuming that transformational
faulting is primarily responsible for deep earthquakes, rather than
alternative mechanisms \citep[e.g.,][]{zhan2020}, our simulations
suggest that coexistence of deep seismicity with continuous slab descent
through the 410 requires a delicate balance: olivine \(\Leftrightarrow\)
wadsleyite reaction rates must be slow enough to sustain metastable
volumes for transformational faulting, yet fast enough to permit 410
penetration.

Rheological contrasts further modulate these kinetic effects by
controlling slab trajectory and descent rate. Strong slabs (\(B\)
\(\gtrsim\) 6, viscosity contrast \textasciitilde6--20\(\times\))
maintain coherence and penetrate sub-horizontally, spending more time
within the phase transition zone where olivine transforms more
completely. Weaker slabs (\(B\) \(\lesssim\) 6, viscosity contrast
\textasciitilde2--6\(\times\)) deform internally and descend vertically,
traversing the transition zone quickly with less time for sluggish
reactions. This morphological dependence on slab strength is consistent
with previous numerical studies without reaction kinetics
\citep{garel2014}, indicating that rheological control of slab geometry
operates independently of kinetic effects.

Considering both kinetic and rheological influences can produce
counterintuitive subduction dynamics, however. Under equivalent kinetic
conditions, our simulations show stronger slabs descending through the
410, while weaker slabs accumulate more metastable olivine, amplifying
buoyancy forces that slow or arrest downward slab motion. This outcome
appears contradictory to previous work linking deep earthquakes with
colder, stronger slabs \citep{green1995, houston2015, zhan2020},
suggesting that other factors---such as overriding plate forcing or
subduction geometry \citep[e.g.,][]{agrusta2017}---may dominate over
kinetics in natural systems.

Our slab simulations also reveal a systematic departure from observed
descent rates across the current global subduction system. Simulated
slabs descend at 1.3--2.0 cm/yr (interquartile range excluding fully
stagnated cases; see Supplementary Information), roughly half of the
range compiled by \citet{lallemand2005} for natural systems
(interquartile range: 2.6--4.8 cm/yr). This discrepancy suggests that
buoyancy forces from metastable olivine accumulation may be exaggerated
in our simulations relative to Earth, although faster descent rates
observed in natural subduction zones should amplify metastability if
kinetic rates are sluggish.

Furthermore, our simulations point to a critical kinetic threshold near
\(\dot{X}\) \(\sim\) 0.07 Ma\(^{-1}\) that marks where buoyancy forces
from incomplete olivine \(\Leftrightarrow\) wadsleyite transformation
either prevent or permit downward slab motion. This narrow threshold
shifts with slab strength, where coherent slabs penetrate the 410 at
reaction rates that would stall weaker slabs (Figure
\ref{fig:410-structure-comp}). These dynamic sensitivities suggest
individual subduction zones could oscillate between penetration and
temporary stagnation at the 410 as kinematic and PT conditions evolve,
with plate forcing and slab deformation behavior
\citep[e.g.,][]{agrusta2017} potentially modulating this critical
kinetic threshold.

Regional variability in slab behavior \citep{fukao2013} could therefore
reflect diversity in both kinetics and rheology, in addition to slab age
and convergence rate. In our simulations, young, hot slabs with lower
viscosity contrasts descend steeply and accumulate moderate metastable
olivine despite warm thermal structures, while old, cold slabs with high
viscosity contrasts flatten and sink slowly, allowing near-complete
transformation under moderately sluggish kinetics. These patterns
distinguish slabs from plumes. In hot upwellings, elevated temperatures
suppress sensitivity to both kinetics and rheology, while in cold slabs,
kinetics govern reaction rates and rheology dictates reaction duration
by regulating slab kinematics. Accurately modeling slab morphology,
material exchange, and deep stress conditions in geodynamic simulations
therefore requires incorporating both effects.

\subsection{Implications for 410 Detectability and Constraining Kinetics
and Rheology from Seismic Observations}\label{sec:410-detectability}

The 410 discontinuity is readily detected using underside reflections
(precursors to SS, PP, and ScS reverberated phases), as well as
converted phases, (receiver functions), and top-side triplicated phases
\citep[e.g.,][]{flanagan1999, shearer2000, chambers2005, schmerr2007, deuss2009, houser2010, cottaar2016, jenkins2016, wang2017, han2021, glasgow2024, miao2024, parera2024, waszek2024}.
The majority of studies focus on mapping 410 topography, and typically
assume a sharp near-horizontal interface. Resolving the width of the
phase transition zone requires more challenging analyses, including
constraints on impedance contrast. The trade-off between phase
transition width, impedance contrast, and small-scale topography depends
strongly on the wavelengths of the seismic phases employed
\citep{chambers2005}.

Previous studies targeting the width of discontinuities have exploited
the frequency dependence of observed amplitudes
\citep[e.g.,][]{vandermeijde2003, bonatto2020}, the detectability of
discontinuities at higher frequencies
\citep[e.g.,][]{rost2002, day2013}, gradient-based inversions
\citep{schmandt2012}, and joint analyses of top- and bottom-side
reflections \citep{vinnik2020}. In many of these approaches, the
transition is parameterized as a simple linear gradient, whereas our
models predict a more complex depth-dependent gradient, consistent with
results from gradient-based inversions \citep{schmandt2012}. While many
of these studies infer a sharp 410, locally broader phase transitions,
up to 35 km, are observed and are commonly attributed to elevated water
content in the mantle
\citep[e.g.,][]{vandermeijde2003, schmandt2012, vinnik2020}.

The presence of water is expected to significantly widen the olivine
\(\Leftrightarrow\) wadsleyite transition width, similar to effects
observed in our simulations. However, a hydrous 410 transition would
likely occur at lower pressures \citep{smyth2002}, while our simulation
results show a deepening of the discontinuity for the intermediate
kinetic regime with an inhibited olivine \(\Leftrightarrow\) wadsleyite
transformation. Therefore, resolving topography could be a means to
distinguish these two scenarios.

Our plume simulations predict consistently thin 410 discontinuities
(2--3 km widths, \textasciitilde24 km displacements) across large
rheological variations and seven orders of magnitude in \(\dot{X}\),
except under ultra-sluggish kinetics where widths broaden up to 21 km.
These sharp discontinuities should be readily detectable, consistent
with observations of well-defined 410s with strong deflections in
various hotspot locations \citep{deuss2009, lawrence2008}.

Slab simulations, in contrast, imply that 410 detectability strongly
depends on joint kinetic and rheological influences (Figure
\ref{fig:410-structure-comp}). Strong slabs (\(B\) \(\gtrsim\) 6,
viscosity contrast \textasciitilde6--20\(\times\)) descending slowly
along sub-horizontal trajectories in our simulations maximize residence
time in the phase transition zone, producing sharp, detectable 410
signals due to near-complete olivine \(\Leftrightarrow\) wadsleyite
transformation. Weak slabs (\(B\) \(\lesssim\) 6, viscosity contrast
\textasciitilde2--6\(\times\)) descending steeply with reduced residence
times amplify kinetic inhibition, producing broader reaction fronts.
Weakened 410 signals beneath some subduction zones
\citep{vanstiphout2019}, or observations of low velocity regions
possibly representing metastable olivine
\citep{jiang2015, shen2020, han2021}, align with intermediate kinetic
conditions and relatively weak slabs, where steep, rapid descent
amplifies metastability.

Understanding the joint influence of rheological and kinetic factors on
the geometry and dynamics of slab descent and phase transition width
will offer a quantitative framework for interpreting observed seismic
structure. Kinetic factors could potentially lead to significantly
extended phase transition zones, but these anomalies are predicted to
occur over narrow regions where slabs interact with the 410 (Figure
\ref{fig:slab-composition-set2}). Thus, finer-scale topographic mapping
with higher frequency phases, such as receiver functions, is likely
needed to localize sharp topography or non-detectability
\citep{vanstiphout2019}.

In principle, regional variations in both 410 topography and phase
transition width could reflect variations in kinetics (through
temperature, composition, water content, or grain size), variations in
rheology (through additional effects like accumulated strain), or
compensating variations in both. Without multiple independent
constraints, seismic observations of 410 structure alone cannot uniquely
separate these effects. Examples for independent constraints may include
the presence of deep seismicity, as well as variations in slab thermal
structure, age, and hydration across subduction zones
\citep{agius2017, vanstiphout2019, schmandt2012}.

Forward modeling of seismically observable 410 width and topography
enables testing of whether specific parameter combinations match
observed 410 structure from receiver functions or precursors. However,
until mineral physics experiments
\citep[e.g.,][]{hosoya2005, kubo2004, perrillat2013, perrillat2016, ledoux2023}
reduce uncertainties and the issue of parameter degeneracy has been
sufficiently mitigated through independent constraints, seismic
observations provide order-of-magnitude estimates of effective
\(\dot{X}\) and relative rheological strength. Nevertheless, the
threshold behavior and scaling relationships from our simulations
demonstrate that comprehensive multi-observation approaches are
essential for constraining micro-scale processes governing phase
transitions and their geodynamic consequences.

\section{Conclusions}\label{sec:conclusions}

The olivine \(\Leftrightarrow\) wadsleyite phase transition and
resulting 410 structure are strongly influenced by the combined effects
of reaction kinetics and rheological strength on flow dynamics. We
quantified these effects by integrating an interface-controlled olivine
\(\Leftrightarrow\) wadsleyite growth model with compressible
simulations of mantle plumes and slabs, systematically exploring kinetic
factors spanning seven orders of magnitude. Each simulation was
evaluated across a large range of viscosity contrasts and 410 structure
was determined after 100 Ma.

Our results reveal fundamentally different responses in hot versus cold
environments. Plumes produce consistently sharp discontinuities (2--3 km
wide) across the entire parameter space, implying that seismic
observations beneath hotspots primarily constrain thermal structure near
thermodynamic equilibrium rather than kinetic or rheological parameters.
Slabs exhibit distinct threshold behavior across three kinetic
regimes---quasi-equilibrium, intermediate, and ultra-sluggish---that are
further modulated by viscosity contrasts controlling slab geometry and
transit time through the phase transition zone.

Widespread slab penetration of the 410 in seismic tomography
\citep{fukao2013} requires effective reaction rates exceeding the
ultra-sluggish kinetic regime (\(\dot{X}\) \(\gtrsim\) 0.07
Ma\(^{-1}\)). This critical threshold shifts systematically with
rheological strength, revealing that modest variations in either
reaction rates or viscosity contrasts can produce substantial diversity
in observed 410 topography. Uniquely constraining kinetic versus
rheological contributions requires combining 410 structural observations
with independent constraints from high-resolution tomography and
kinematic data (providing descent geometry) and deep seismicity patterns
(indicating metastable olivine volumes).

The 410 can therefore serve as a seismological probe of kinetic
conditions in cold subduction environments where disequilibrium effects
are amplified. Realizing this potential requires determining
discontinuity width through targeted seismic studies as well as reducing
uncertainties through targeted mineral physics experiments that better
quantify nucleation versus growth mechanisms, water and compositional
effects on reaction rates, and microstructural evolution during
deformation. Integrating such constraints with high-resolution seismic
imaging and the forward modeling framework presented here offers a
pathway toward understanding how micro-scale kinetic processes govern
mantle convection and shape Earth's interior seismic structure.

\section*{Acknowledgements}\label{acknowledgements}
\addcontentsline{toc}{section}{Acknowledgements}

This work was funded by the UKRI NERC Large Grant no. NE/V018477/1. All
computations were undertaken on Barkla2, part of the High Performance
Computing facilities at the University of Liverpool, who graciously
provided expert support. We thank the Computational Infrastructure for
Geodynamics (\url{https://geodynamics.org}) which is funded by the
National Science Foundation under award EAR-0949446 and EAR-1550901 for
supporting the development of ASPECT.

\section*{Data Availability}\label{sec:data-availability}
\addcontentsline{toc}{section}{Data Availability}

All data, code, and relevant information for reproducing this work can
be found at
\url{https://github.com/buchanankerswell/kerswell_et_al_410_kinetics},
and at \url{https://doi.org/10.17605/OSF.IO/9PHWC}, the official Open
Science Framework data repository. All code within these repositories is
MIT Licensed and free for use and distribution (see license details).
ASPECT version 3.0.0,
\citep{aspect-doi-v3.0.0, aspectmanual, heister2017, kronbichler2012, gassmoller2018, clevenger2021, fraters2019, fraters2020}
used in these computations is freely available under the GPL v2.0 or
later license through its software landing page
\url{https://geodynamics.org/resources/aspect} or
\url{https://aspect.geodynamics.org} and is being actively developed on
GitHub and can be accessed via
\url{https://github.com/geodynamics/aspect}.

\section*{Conflict of Interest}\label{sec:conflict-of-interest}
\addcontentsline{toc}{section}{Conflict of Interest}

The authors declare there are no conflicts of interest for this
manuscript.

\clearpage

\bibliography{/Users/bck/Working/manuscripts/kerswell_et_al_410_kinetics/draft/main.bib}


\end{document}
